\section{Automatic induction and demand-driven lemma discovery}
\status{The list fragment, changing-argument detection, bounded RHS synthesis,
structural induction, and trace replay are implemented. General dependent
induction and e-graph-guided frontier discovery are specified but not implemented.}

\subsection{Choose an induction variable from recursive structure}

Collect recursive applications in the definitions mentioned by the goal and
its relevant hypotheses. For each function, record which argument positions
are destructed by pattern matching, which are passed to recursive calls as
proper subterms, and which non-structural arguments change. Rank induction
candidates by how many blocked recursive calls they expose after one
constructor split. Do not select a variable merely because its type has an
induction principle.

For a first implementation, restrict automatic induction to standard
structural recursion over supported inductive types. A later adapter can use
functional or well-founded induction principles associated with a recursive
definition. It must construct an actual Lean recursor application; recognizing
a decreasing-looking expression is not a replacement for a termination or
well-foundedness proof.

The candidate motive is a typed function from the induction index to a
proposition. Generalize parameters whose values vary in recursive calls.
Close that set under dependencies in the local telescope: if a generalized
variable occurs in the type of another variable or hypothesis, that dependent
binder must be generalized or explicitly transported as well.

For instance, in a context containing $x:A$, $y:B(x)$, and $h:C(x,y)$,
induction on $x$ may require a motive of the form
\[
M(x)\;:=\;\forall y:B(x),\; C(x,y)\to P(x,y),
\]
not a proposition containing the old fixed $y$ at every new value of $x$.
In Lean, use elaboration and the existing induction machinery to build this
motive; do not manufacture a well-typed-looking first-order erasure.

\subsection{The accumulator example, including the obstruction}

Write $A$ for append, $R$ for ordinary reverse, and $\operatorname{RA}$ for
accumulator reverse. The defining equations are
\begin{align*}
A([],y)&=y,& A(h::t,y)&=h::A(t,y),\\
R([])&=[],&R(h::t)&=A(R(t),[h]),\\
\operatorname{RA}([],a)&=a,&
\operatorname{RA}(h::t,a)&=\operatorname{RA}(t,h::a).
\end{align*}
Suppose the requested theorem is
\[
\operatorname{RA}(x,[])=R(x).
\]
Inducting on $x$ without generalization gives an induction hypothesis only
about $\operatorname{RA}(t,[])$. But the recursive call in the step case is
$\operatorname{RA}(t,[h])$. The induction hypothesis does not match.
Increasing a rewriting budget cannot make this particular hypothesis say
something stronger.

The changing-argument detector sees that argument 0 is destructed structurally
and argument 1 changes from $a$ to $h::a$. It replaces the fixed empty
accumulator with a fresh parameter and asks for a right-hand side in a typed
grammar:
\[
L ::= []\mid x\mid a\mid R(L)\mid A(L,L).
\]
The successful candidate is
\begin{equation}\label{eq:accspec}
\forall x\,a,\qquad \operatorname{RA}(x,a)=A(R(x),a).
\end{equation}
The base case is immediate. For the step case, the generalized induction
hypothesis gives
\[
\operatorname{RA}(t,h::a)=A(R(t),h::a).
\]
The target's other side is
\[
A(R(h::t),a)=A(A(R(t),[h]),a)
          =A(R(t),A([h],a))=A(R(t),h::a),
\]
where the middle equality uses append associativity. Instantiating
\eqref{eq:accspec} at $a=[]$ and using append's right identity proves the
original theorem.

The implemented grammar enumerator finds this RHS as candidate 32. It filters
candidates using all 49 pairs of lists of length at most two over the alphabet
$\{0,1\}$. Exactly one candidate reaches the proof search before success.
These finite examples are \emph{not} its proof: the candidate then passes
symbolic induction and independent trace replay.

\subsection{What the miniature checker actually permits}

There are separate syntax constructors for universal pattern variables and
rigid free variables. In an induction step, the tail is a rigid free variable.
Other generalized parameters remain variables of the induction-hypothesis
rule. Thus the hypothesis can be instantiated at a new accumulator, but cannot
be applied to an arbitrary longer list.

Every rewrite certificate contains a rule identifier and a subterm path. The
checker retrieves the rule from the defining equations, an earlier checked
lemma, or the current local induction hypothesis. It matches the left-hand
side, checks repeated-variable identity and sorts, replaces the selected
subterm, and continues. The endpoints of the two replayed sides must be
syntactically identical.

The public theorem-installation operation first checks both induction cases.
A lemma is not available while its own certificate is being checked. This
rejects a forged proof that cites its own conclusion. The prototype does not
implement simultaneous induction; a production extension would need an
explicit joint motive and recursor rather than relaxing this restriction.

\begin{proposition}[Soundness of the list certificate format]
Assume that previously installed equations are valid for all finite lists.
If the miniature checker accepts an induction certificate, its universally
quantified equation holds for all finite lists under the supplied definitions.
\end{proposition}
\begin{proof}
Each defining rewrite preserves the interpreted value. Each installed lemma
instance preserves it by the assumption, and congruence permits replacement
inside a typed term context. Consequently an accepted base trace proves the
empty-list case. In the step trace, the only additional rule is the desired
statement for the rigid tail, universally quantified over its generalized
parameters. That rule is valid under the structural induction hypothesis.
The accepted step trace therefore proves the constructor case. Structural
induction establishes the universal equation.
\end{proof}

This is a mathematical justification of a small format, not a machine-checked
Lean proof of the Python implementation. Bugs in the checker remain possible;
the negative tests and replay separation reduce that risk but do not turn
Python into Lean's kernel.

\subsection{Discover the helper lemma, not only the generalization}

The prototype supplies four seed \emph{statements}: append associativity,
append's right identity, reverse of append, and reverse involution. It searches
and checks their proofs, but does not claim to have invented all four
statements. In particular, associativity is available when the accumulator
lemma is synthesized.

The full design adds a demand-driven lemma explorer. After a failed induction
step, retain the normalized left and right frontiers and the rewrites that
led to them. Bounded equality saturation can expose alternative forms that a
one-way normalizer would discard. For the accumulator example, unfolding
append in reverse can expose $h::a$ as $A([h],a)$. The remaining frontier then
suggests
\[
A(A(u,v),w)=A(u,A(v,w)).
\]
Prove that generalized statement separately, and only then install it.

A practical candidate-generation algorithm is:
\begin{lstlisting}
find_lemma(left_frontier, right_frontier):
  expand each frontier using a bounded set of proved equalities
  identify pairs separated by a small structural difference
  anti-unify shared subterms, preserving repeated-variable identity
  abstract a well-typed parameter telescope
  reject trivial, cyclic, or irrelevant candidates
  filter candidates with executable counterexamples where available
  prove surviving candidates in a separate obligation
  return only checked lemmas to the main equality state
\end{lstlisting}

Typed anti-unification starts in a deliberately restricted fragment:
fully applied, homogeneous first-order terms with compatible types. Memoize
pairs of subterms so that anti-unifying $f(a,a)$ and $f(b,b)$ gives $f(z,z)$,
not $f(z_1,z_2)$. In a dependent fragment, any abstraction that changes a term's
type needs a typed telescope and possibly transport proofs. Failure to find
such an abstraction is an ordinary search failure.

This goal-directed exploration is strongly related to CCLemma's approach.
Its authors use e-graphs to focus inductive lemma discovery on useful proof
frontiers. The present proposal adapts that line of work to a Lean
proof-producing controller; it does not claim to originate the method.
\cite{cclemma}

\subsection{Budgeting and termination of exploration}

Impose separate limits on term size, saturation nodes, candidate count,
induction depth, and lemma-dependency depth. Use term fingerprints and
alpha-normalized candidate telescopes to avoid reproving the same lemma under
new binder names. Never accept a cyclic dependency merely because every edge
looks like a useful lemma application.

Bounded enumeration is exponential in its grammar size. The initial list
experiment is intentionally tiny. The production advantage must come from
recursive-call analysis, relevance filtering, typed indexing, and frontier
selection, not from assuming unrestricted enumeration will be cheap.

\section{Polynomial recurrence invariants by exact CEGIS}
\status{Implemented and tested on 65 recurrence problems, including five power-sum
families and 60 seeded random polynomial increments.}

\subsection{The fragment and the certificate}

Consider the recurrence
\begin{equation}\label{eq:recurrence}
T(0)=t_0,\qquad T(n+1)=T(n)+f(n),\qquad n\in\N,
\end{equation}
where $t_0\in\Q$ and $f\in\Q[X]$. Search for $q\in\Q[X]$ such that
\begin{equation}\label{eq:rec-cert}
q(0)=t_0,\qquad q(X+1)-q(X)=f(X).
\end{equation}
The certificate is just the exact coefficient list of $q$. Its checker
computes the shift $q(X+1)$ using polynomial arithmetic and compares exact
coefficients. It does not sample the step relation or trust an interpolator.

\begin{proposition}
If \eqref{eq:rec-cert} holds, then $T(n)=q(n)$ for every $n\in\N$.
\end{proposition}
\begin{proof}
The base equality is the first certificate condition. If $T(n)=q(n)$, then
\eqref{eq:recurrence} and the second condition imply
$T(n+1)=q(n)+f(n)=q(n+1)$. Apply induction on $n$.
\end{proof}

In Lean, this certificate can be reconstructed with ordinary induction,
explicit casts, and polynomial normalization. A compact reflected checker is
an optimization, not a requirement for the first implementation.

\subsection{Counterexamples direct the coefficient search}

For degree bound $d$, write $q(X)=\sum_{i=0}^d c_iX^i$. Maintain exact sample
constraints $q(n)=T(n)$ and solve the resulting rational linear system.
Free coefficients are initially set to zero. Check the candidate's base and
step identities. If the step residual
\[
r(X)=q(X+1)-q(X)-f(X)
\]
is nonzero, choose a natural $n$ with $r(n)\ne0$ and add both $q(n)=T(n)$ and
$q(n+1)=T(n+1)$ to the sample set. At least one is a new constraint: if both
were already satisfied by the candidate, their difference would force
$r(n)=0$.

A nonzero degree-$k$ polynomial cannot vanish at all of $0,1,\ldots,k$.
The implementation uses this elementary fact to find an exact counterexample
without an SMT solver. If the coefficient constraints become inconsistent,
increase the degree bound. The proposal is therefore counterexample-guided,
while its acceptance test is an exact induction certificate.

\subsection{A completeness statement for this small fragment}

For a nonzero polynomial $f$ of degree $k$, the finite-difference map
$q\mapsto q(X+1)-q(X)$ sends a leading term $aX^{k+1}$ to
$(k+1)aX^k$ plus lower terms. Over $\Q$, solve for the leading coefficient,
subtract its difference, and continue downward. Thus there is a degree
$k+1$ solution, unique after fixing $q(0)$. The zero increment gives a constant
solution.

For a fixed candidate degree, each unsuccessful iteration adds an independent
constraint or detects inconsistency. Hence exact linear solving and a degree
budget at least $k+1$ suffice for termination on this fragment. This does not
assert completeness for arbitrary recurrences, nonlinear state updates,
piecewise definitions, or integer formulas involving truncated division.

\subsection{An example and an overfitting trap}

For $f(X)=X+1$, the discovered result is
\[
q(X)=\frac{X(X+1)}2.
\]
For increments $(X+1)^2,\ldots,(X+1)^5$, the output contains the corresponding
power-sum polynomials; the exact formulas and all failed candidate residuals
are in \code{results/certificates/recurrences.json}.

An interpolated expression
\[
\widetilde q(X)=\frac{X(X+1)}2+X(X-1)(X-2)(X-3)
\]
agrees with the triangular-number formula on the first four inputs. It is
nevertheless rejected by the step checker. This is a small, concrete example
of why testing-based conjecture generation and proof acceptance must be
separate operations.

\subsection{Generalization to richer invariants}

A next worker can search polynomial invariants $I(n,s)$ for a state transition
$s'=F(n,s)$. It must prove the initial condition and the implication
\[
I(n,s)\land \operatorname{Guard}(n,s)
\;\Longrightarrow\; I(n+1,F(n,s)).
\]
Coefficient comparison handles identities; exact nonlinear certificates can
handle some inequalities. The invariant template and degree bound are explicit
search parameters. Counterexamples refine coefficients or force a richer
motive. This is a concrete extension path, but the supplied recurrence module
implements only the additive univariate fragment above.

\section{Synthesize witnesses, not merely satisfying assignments}
\status{Affine witness synthesis with exact Farkas certificates is implemented.
Piecewise witnesses, integer witnesses, and dependent-type witness synthesis
are proposed extensions.}

\subsection{The logical distinction}

To prove $\forall x\,\exists y\,P(x,y)$, a satisfying assignment for one value
of $x$ is not enough. A useful synthesizer produces a term $y=w(x)$ and proves
$\forall x\,P(x,w(x))$. For example,
\[
\forall x\in\R,\quad \exists y\in\R,\quad x+1\le y\le x+2
\]
has the witness $w(x)=x+1$. No constant witness works for every real $x$.
The prototype's constant-only template returns no certificate; its affine
template returns the source-level witness $x+1$ with an exact check.

In the Lean controller, witness binders retain their dependencies. A witness
for the second existential may refer to an earlier witness, but not to a
future binder or a branch-local variable outside its scope. A witness is
elaborated as an ordinary Lean term before its specification is discharged.

\subsection{Jointly solve for an affine witness and its proof}

Let the input domain be the rational polyhedron
\[
g(x)=Ax+b\ge0,
\]
and the target constraints be
\[
h(x,y)=Cx+Dy+e\ge0.
\]
All inequalities are componentwise and non-strict. Seek $y=Mx+c$. A sufficient
certificate consists of a nonnegative matrix $\Lambda$ and nonnegative vector
$s$ satisfying
\begin{align}
C+DM&=\Lambda A,\label{eq:farkas-linear}\\
e+Dc&=\Lambda b+s.\label{eq:farkas-constant}
\end{align}
These equations are \emph{linear in the unknowns} $M,c,\Lambda,s$, because
$A,b,C,D,e$ are fixed problem data. Witness synthesis and proof search can
therefore be one linear feasibility problem, rather than enumerating many
candidate affine functions and invoking a verifier separately.

\begin{proposition}[Affine witness certificate]
If \eqref{eq:farkas-linear}--\eqref{eq:farkas-constant} hold and
$\Lambda,s\ge0$, then for every $x$ satisfying $Ax+b\ge0$, the term
$y=Mx+c$ satisfies $Cx+Dy+e\ge0$.
\end{proposition}
\begin{proof}
Substitution gives
$Cx+D(Mx+c)+e=\Lambda(Ax+b)+s$.
Every entry on the right is a sum of nonnegative multiples of nonnegative
quantities, plus a nonnegative slack.
\end{proof}

For a nonempty rational polyhedron, the affine implication form of Farkas'
lemma also gives a completeness route for a \emph{fixed affine witness}:
validity of each non-strict affine target inequality has such multipliers.
The combined exact feasibility formulation therefore captures the existence
of affine witnesses in that fragment. This is the standard polyhedral dual-certificate reasoning.\cite{mosek}
The floating-point proposer is not a complete exact LP algorithm, and an empty domain needs separate handling
for a completeness claim. Neither qualification affects certificate soundness.

\subsection{Implementation details and exact reconstruction}

With $n$ inputs, $m$ outputs, $r$ domain constraints, and $q$ target constraints,
the encoding uses
\[
m(n+1)+q(r+1)
\]
unknowns and $q(n+1)$ coefficient equations, before template restrictions.
Entries of $M,c$ are unrestricted; entries of $\Lambda,s$ are constrained
nonnegative. A constant-only grammar adds equations $M=0$.

The prototype calls SciPy's \code{linprog} with the HiGHS backend. That API
is an optimization interface, not a proof system.\cite{scipylp} The returned
numbers are rationalized and all coefficient equations and sign constraints
are checked with exact fractions. If that fails, a second attempt uses the
numerically suggested support and exact rational elimination. If the exact
check still fails, the answer is ``no certificate found.''

A numerical zero threshold is used only to select a support for the second
search attempt. There is no numerical tolerance in the acceptance condition.
The serialized certificate includes each witness, each multiplier row, and
each slack. Its replay code imports no optimization library.

For Lean reconstruction, produce the witness tuple first, then prove every
target inequality from the recorded nonnegative combination. The coefficient
identity can be checked by polynomial normalization, and the sign proof can
be built from the domain hypotheses. One can use ordinary tactic reconstruction
initially and a reflected matrix checker later.

\subsection{Beyond a single affine template}

There are valid specifications with no witness in the chosen grammar.
For piecewise-affine witnesses, permit a bounded decision tree whose guards
are source-level affine comparisons. Synthesize an affine witness under each
leaf's domain assumptions, and prove the guard partition covers the original
domain. Never silently assume that finitely many sampled regions cover it.

For integer or natural-number witnesses, a rational affine solution is not
sufficient. Require integral coefficients where that suffices, or introduce
explicit division/remainder templates with associated arithmetic obligations.
A practical first integer grammar is
\[
w(x) ::= c+\sum_i a_ix_i\mid
\left\lfloor\frac{c+\sum_i a_ix_i}{d}\right\rfloor\mid
\operatorname{ite}(g,w_1,w_2),\quad d>0.
\]
Each candidate must be checked in the source integer semantics. The current
prototype deliberately does not interpret a real solution as an integer
proof.

\section{Theory-constrained quantifier instantiation}
\status{Specified algorithm; no general E-matching or MBQI implementation is included
in the prototype archive.}

\subsection{Solve for missing terms before increasing instance limits}

Trigger-based instantiation works from terms present in the equality state.
The documented E-matching machinery is therefore an important baseline, not
something to discard.\cite{grindematch} \Forge{} adds a demand that works
backward from a desired consequence to source terms for the theorem's binders.

Consider an integer predicate $P$, hypotheses
\[
\forall k\in\Z,\ P(2k+1),\qquad n\bmod2=1,
\]
and goal $P(n)$. A useful instantiation is $k=n\mathbin{/}2$. The controller
can synthesize that term from the arithmetic shape, then prove
\[
2(n\mathbin{/}2)+1=n
\]
using the oddness assumption and integer division facts. It applies the
original universal theorem and transports along this checked equality.

This is an algorithmic example, not a measured claim that the current
\Grind{} fails on that exact input. The included Lean example shows the
intended proof shape but was not compiled.

\subsection{A concrete constrained-matching algorithm}

Given a selected theorem conclusion $C(z)$ and a requested atom $G$, first
perform typed syntactic matching on the uninterpreted structure. Instead of
failing at an interpreted arithmetic subterm, emit an equality constraint
between that subterm and the corresponding source expression.

For linear integer patterns, solve the resulting constraints for missing
binders using divisibility, linear elimination, or a bounded grammar of
quotient/remainder terms. For affine rational patterns, exact linear solving
can directly produce terms. For constructor patterns, use constructor
injectivity and source-level constructor terms. Underconstrained binders are
handled by a bounded typed grammar, not arbitrary model objects.

Each candidate substitution $\theta$ creates ordinary proof obligations for
the matching equalities and any theorem premises. After those are proved,
construct the theorem application $C(\theta)$ and its transport to $G$.
Nothing about the search solver needs to be trusted: an incorrect candidate
simply fails reconstruction.

Several triggers or repeated binders produce a \emph{joint} constraint system.
Solving each occurrence independently would lose the equality relationship
between two occurrences of the same variable. Cache instances using the
theorem identity, universe/type instantiations, the complete substitution,
and its valid branch context.

\subsection{Model-guided instantiation, with a narrow trust boundary}

If constrained matching is insufficient, use a model-based loop inspired by
quantified SMT procedures. Candidate-model and pattern-based instantiation
are established techniques in systems such as Z3.\cite{z3quant} They can
suggest useful instances even when naive trigger firing stagnates.

The proposed loop freezes a relevant typed abstraction, obtains a candidate
model, and searches for a violated universal instance in that model. It then
tries to reify the candidate values as source terms and verifies the associated
instance in Lean. Numerals can often be reified directly. An uninterpreted
model-domain element cannot simply become a new arbitrary Lean constant:
find a denoting term already available, enumerate a term, or abandon that
candidate. Internal Skolem symbols likewise require scoped witnesses or
hypothesis elimination, not new axioms.

A satisfying model of an abstraction is not automatically a countermodel of
the Lean goal. Conversely, an external unsatisfiability result is not accepted
without proof reconstruction or an appropriate verified certificate. In this
architecture, model search is an instance generator; Lean remains responsible
for every asserted consequence.

Generation limits, trigger selectivity, repeated-instance suppression, and
fair round scheduling are still necessary. There is no claim that this loop
is complete for arbitrary quantified dependent type theory.

\section{Nonlinear reasoning with exact certificates}
\status{Finite-dictionary square search and Bernstein subdivision are implemented.
A full SDP solver and its Lean integration are reused/proposed, not implemented
by these Python modules.}

\subsection{A nonlinear worker should expose multiple certificate classes}

Start with cheap normalization and obvious sign facts. Use Mathlib's
\code{positivity} and arithmetic tactics as proof-producing preprocessors or
closers where appropriate.\cite{positivity,linarith} Then select a more
specialized certificate class according to the target and assumptions:
algebraic identities, nonnegative square combinations, bounded-box Bernstein
certificates, or an optional stronger semialgebraic backend.

Do not confuse polynomial equality reasoning with order reasoning. Reducing
polynomials modulo equalities does not by itself prove nonnegativity.
Conversely, knowing a square is nonnegative does not solve all nonlinear
integer arithmetic problems.

\subsection{Square certificates and the existing Lean SOS backend}

For constraints $g_i(x)\ge0$ and equations $e_j(x)=0$, a useful certificate
shape is
\begin{equation}\label{eq:sos}
p=\sum_k a_kq_k^2
   +\sum_{i,k} b_{ik}r_{ik}^2g_i
   +\sum_j h_je_j,
\qquad a_k,b_{ik}\in\Q_{\ge0}.
\end{equation}
Checking the polynomial identity and the rational signs proves $p\ge0$ under
the source constraints. The polynomial multipliers $h_j$ of equalities need
no sign condition.

There is already a Lean \code{sos} project implementing proof-producing
nonlinear reasoning with rational certificates. Its documented architecture
separates reification, search, checking, and verification; named interfaces
include \code{SOS.Engine.solve}, \code{SOS.Engine.Result.check}, and the
Mathlib-side verifier. \code{sos?} can emit a replayable witness. Its
computational polynomial substrate uses \code{Hex.MvPoly}.\cite{sos}
These are concrete integration targets, subject to revision compatibility.
The \Forge{} contribution is when to request such a certificate and how to
return its proved consequences to structural and saturation workers.

\subsection{What the implemented square search does}

The prototype uses only the first two sums of \eqref{eq:sos}. It chooses a
finite dictionary of candidate square roots, constructs their squares and
constraint-multiplied squares, and solves for nonnegative rational weights.
Once the roots are fixed, coefficient matching is a linear feasibility
problem. For two variables and root degree at most two, the default dictionary
has 36 roots: monomials and their pairwise sums and differences.

This is simpler and much less general than searching over arbitrary positive
semidefinite Gram matrices. It is useful as a cheap stage and as an experiment
in exact search/replay separation, not as a proposed replacement for a full
SOS implementation. Its failure proves only that this search did not find a
certificate. In particular, a nonnegative polynomial need not be a sum of
polynomial squares at all.

The checker reconstructs the whole polynomial from the supplied weights and
roots, checks each weight exactly, verifies constraint indices, and compares
the result with the target supplied by the caller. A weight of $-1$ cannot be
smuggled through because the polynomial identity happens to be correct.
Removing a necessary constraint from the problem also makes replay fail.

\subsection{Bernstein certificates on rational boxes}

Suppose the domain is a rational box
$B=\prod_{i=1}^m[a_i,b_i]$, with $a_i<b_i$. Substitute
$x_i=a_i+(b_i-a_i)t_i$ to obtain
\[
\widetilde p(t)=\sum_\beta c_\beta t^\beta,\qquad 0\le t_i\le1.
\]
Let $d_i$ be a coordinatewise degree bound. Define
\[
B_{\alpha,d}(t)=\prod_i\binom{d_i}{\alpha_i}
 t_i^{\alpha_i}(1-t_i)^{d_i-\alpha_i}.
\]
Each basis function is nonnegative on the unit box and their sum is one.
The Bernstein coefficients are
\begin{equation}\label{eq:bernstein}
b_\alpha=\sum_{\beta\le\alpha}c_\beta
 \prod_i\frac{\binom{\alpha_i}{\beta_i}}{\binom{d_i}{\beta_i}}.
\end{equation}
Thus $p\ge0$ on $B$ if every $b_\alpha\ge0$; it is strictly positive if every
$b_\alpha>0$.

\begin{proposition}[Bernstein coefficient identity]
The coefficients in \eqref{eq:bernstein} satisfy
$\widetilde p(t)=\sum_\alpha b_\alpha B_{\alpha,d}(t)$.
\end{proposition}
\begin{proof}
In one variable, use
\[
 t^k=\sum_{j=k}^d
 \frac{\binom{j}{k}}{\binom{d}{k}}
 \binom{d}{j}t^j(1-t)^{d-j}.
\]
The binomial identity
$\binom{d}{j}\binom{j}{k}/\binom{d}{k}=\binom{d-k}{j-k}$
reduces the right-hand side to
$t^k(t+(1-t))^{d-k}=t^k$. Take products across coordinates and apply linearity
to the monomial expansion of $\widetilde p$.
\end{proof}

Negative Bernstein coefficients are \emph{not} counterexamples. They mean
only that this sufficient condition has not proved the desired sign. A
search worker can split a box at a rational midpoint and repeat the check.
A returned negative point is a different result and must be evaluated exactly.

\subsection{The partition certificate is part of the proof}

The implemented certificate is a binary tree. A leaf asks the checker to
recompute the Bernstein coefficients for its inherited box. An internal node
specifies an axis, a rational split point strictly inside that axis's
interval, and both children. The child boxes are derived by the checker, not
accepted as arbitrary boxes from the certificate.

This design makes coverage explicit. There is no way to omit a difficult
corner by supplying an incomplete list of favorable boxes. A missing child,
an invalid split axis, or a split point at an interval endpoint is rejected.
The root box and polynomial are separate problem inputs.

For
\[
p(x)=\left(x-\tfrac12\right)^2+\tfrac1{100},\qquad x\in[0,1],
\]
the degree-two coefficients at the root are
$13/50,-6/25,13/50$. The root check cannot certify positivity. Splitting at
$1/2$ gives coefficient lists
\[
(13/50,1/100,1/100),\qquad(1/100,1/100,13/50),
\]
so two leaves suffice for strict positivity. This is an exact calculation,
not a numerical plot or an assertion based on sampled values.

\subsection{Limits, completeness, and complementary strengths}

For a polynomial strictly positive on a compact box, sufficiently fine
subdivision eventually makes every Bernstein coefficient positive. A direct
justification is that the transformed higher-degree coefficients tend
uniformly to zero as box widths shrink, while the constant terms stay above
a positive minimum. Consequently every local Bernstein coefficient is
$p(a)+O(\max_i(b_i-a_i))$ with a uniform bound over the original box.
The prototype cycles split axes so every coordinate shrinks along a branch.
An unlimited fair subdivision search therefore has a useful completeness
property for strictly positive polynomials on a fixed compact box.

Nonnegative polynomials with zeros are different. The prototype returns
\code{unknown} for $(x-y)^2$ on $[0,1]^2$ at depth four, even though a single
square certificate proves the result immediately. A search timeout or an
uncertified leaf must not be reported as a refutation.

Coordinatewise degrees create $\prod_i(d_i+1)$ Bernstein coefficients, and a
binary tree can have exponentially many leaves in its depth bound. Exact
coefficient bit sizes also grow under repeated substitution. This procedure
is attractive for low-dimensional bounded problems, not a universal substitute
for nonlinear arithmetic.

\subsection{Integrating domains, casts, and division safely}

The controller must establish that the source assumptions imply the box bounds
before applying a Bernstein theorem. Polynomial reification must preserve
variable identity, source operations, and casts. A real-field certificate
cannot be applied directly to an arbitrary ring or a finite characteristic
field.

For natural or integer inequalities, explicitly prove the embedding into an
ordered field and translate the relevant operations. Integer division,
truncated subtraction, and modular arithmetic are not field operations.
Clearing denominators in inequalities requires a sign proof; clearing a field
denominator in an equality generally requires a nonzero proof. These become
\code{ProveGuard} obligations. If the guard fails, the rewrite is unavailable,
not approximately valid.

A checked sign or equality from this worker is useful beyond its original
goal. It can unlock a conditional rewrite, expose a useful quantified
instance, or finish an induction step. That is why the arithmetic worker
returns individual proved consequences and their assumptions, not just a
Boolean ``solved'' flag.
